% ── SECTION 4: THE CHRONOLOGICAL SEQUENCE (~700 words) ──────────────────────

\section{The Chronological Sequence: What It Shows}
\label{sec:chronology}

% Present convergence table in narrative form.
% What the sequence shows and what it does not prove.

The convergence I have documented has a chronological dimension
that is part of my argument.
A finding of structural identity between two traditions might, in principle,
be explained by influence: one tradition learns the structure from the other
and reformulates it in its own vocabulary.
The chronology shown in Table~\ref{tab:chronology} forecloses that explanation.
The \bahai articulations precede the philosophical and physical articulations
by margins that make influence impossible in either direction.
(The characterizations of the \bahai writings in the table reflect my analytical reading;
the texts themselves did not employ the vocabulary of process ontology or relational QM.)

\begin{table}[h]
\centering
\caption{Chronological sequence of convergent findings.}
\label{tab:chronology}
\begin{tabularx}{\textwidth}{llX}
\toprule
\textbf{When} & \textbf{By Whom / How} & \textbf{What Was Established} \\
\midrule
1850s--1892  & \Bahaullah (revelation)
             & Process ontological claims: reality as motion, creation as
               relational event, existence as relational affinity \\
1892--1921   & \AbdulBaha (teaching)
             & Love as creative foundation; hierarchy of comprehension;
               soul as non-composite \\
1929         & Whitehead (philosophical analysis)
             & \textit{Process and Reality}: actual occasions, prehension,
               eros --- process ontology named philosophically \\
1964         & Bell (mathematics)
             & Bell's theorem: no local hidden-variable theory reproduces quantum predictions \\
1982         & Aspect, Grangier, Roger (experiment)
             & Landmark test: Bell inequalities violated, as quantum mechanics predicts \\
1996         & Rovelli (theoretical physics)
             & Relational QM: quantum states are always relative, never absolute \\
2022         & Zeilinger, Aspect, Clauser (Nobel Prize)
             & Nobel Prize for experiments establishing the violation of Bell inequalities \\
\bottomrule
\end{tabularx}
\end{table}

The \bahai writings were not responding to quantum mechanics --- they
preceded it by a century.
They were not responding to Whitehead --- they preceded him by decades.
They arrived where they arrived by their own means: spiritual discernment,
articulated in the writings of \Bahaullah and \AbdulBaha.

This is not presented as proof of supernatural authority.
What it establishes, on my reading, is independence: three traditions that did not communicate
with each other, using three different methods, arrived at what I read as the same structural
description of reality.
That convergence requires explanation.

A reader may reply that structural coincidences across traditions are not rare
in the history of ideas, and that no particular convergence demands a
particular metaphysics.
The reply is fair as a general caution.
What it does not address is the specificity of the present convergence.
Six independent structural points, drawn from texts that did not communicate
with each other, agreeing not in mood or imagery but in logical architecture:
this is a degree of correspondence for which the appeal to coincidence
becomes itself a load-bearing claim, and one that requires its own argument.
I do not insist on a metaphysical conclusion here.
I insist only that the convergence is real, that it is specific, and that
explanations of it must take its specificity into account.
